% Toolchain test document; will be replaced by the thesis template in the next step.
\documentclass[a4paper, 12pt]{scrreprt}

\usepackage[T1]{fontenc}
\usepackage{lmodern}
\usepackage[ngerman]{babel}
\usepackage{minted}
\usepackage[autostyle=true]{csquotes}
\usepackage{microtype}

\usepackage{tikz}
\usepackage{pgfplots}
\pgfplotsset{compat=1.18}

\usepackage[backend=biber, style=authoryear]{biblatex}
\addbibresource{references.bib}

\usepackage[hidelinks]{hyperref}
\usepackage[acronym, toc]{glossaries-extra}
\setabbreviationstyle[acronym]{long-short}
\makeglossaries

\usepackage[ngerman]{cleveref}
\crefname{listing}{Listing}{Listings}

\newacronym{api}{API}{Application Programming Interface}
\newglossaryentry{latex}{
  name={\LaTeX},
  sort={LaTeX},
  description={Textsatzsystem auf Basis von \TeX}
}

\title{Toolchain-Test}
\author{Vorname Nachname}
\date{\today}

\begin{document}

\maketitle
\tableofcontents

\chapter{Funktionstest}

Umlaute und Silbentrennung: Übergrößenträger, Straßenbahnhaltestelle,
\enquote{deutsche Anführungszeichen}.

Literatur: \textcite{knuth1984} beschreibt \TeX{}, \LaTeX{} wurde von
\citeauthor{lamport1994} entwickelt \autocite{lamport1994}.

Glossar: Eine \gls{api} wird beim zweiten Mal nur als \gls{api} abgekürzt.
Gesetzt wird mit \gls{latex}.

Querverweise: siehe \cref{fig:plot} und \cref{lst:hallo}.

\begin{listing}[htbp]
  \begin{minted}{python}
def gruss(name: str) -> str:
    return f"Hallo, {name}!"
  \end{minted}
  \caption{Python-Beispiel mit minted}
  \label{lst:hallo}
\end{listing}

Inline-Code: \mintinline{python}{print(gruss("Welt"))}.

\begin{figure}[htbp]
  \centering
  \begin{tikzpicture}
    \begin{axis}[width=0.7\textwidth, xlabel={$x$}, ylabel={$x^2$}, grid=major]
      \addplot[domain=-2:2, samples=50, thick, blue] {x^2};
    \end{axis}
  \end{tikzpicture}
  \caption{Testplot mit pgfplots}
  \label{fig:plot}
\end{figure}

\printglossary[type=\acronymtype, title={Abkürzungsverzeichnis}]
\printglossary[title={Glossar}]
\printbibliography

\end{document}
